%------------------------------------------------------------------------------%

\usepackage{integracao_e_series}

\title{Integração por Substituição Trigonométrica}
\subtitle{1 -- Substituição pelo seno}

%------------------------------------------------------------------------------%
\begin{document}

\addtitle

%------------------------------------------------------------------------------%
\section{Integração por Substituição Trigonométrica}

%------------------------------------------------------------------------------%
\begin{frame}{Substituição Simples}
  Calculando
  \[
    F = \int x \sqrt{a^2-x^2}\,dx
  \]
  \pause
  substituição
  \(
    \qquad
    u = a^2-x^2
    \qquad\qquad
    du = -2xdx
    \qquad
    xdx = -\frac{du}{2}
  \)
  \pause
  \[
    F = -\frac{1}{2}\int \sqrt{u}\,du
  \]
\end{frame}

%------------------------------------------------------------------------------%
\begin{frame}{Substituição Simples -- Problema}
  Calculando
  \[
    F = \int \sqrt{a^2-x^2}\,dx
  \]
  \pause
  substituição
  \(
    \qquad
    u = a^2-x^2
    \qquad\qquad
    du = -2xdx
    \qquad
    dx = -\frac{du}{2x}
  \)
  \pause
  \[
    F
    = -\frac{1}{2}\int \sqrt{u} \,\frac{du}{{\color<4->{red}x}}
    \pause
    =\; {\color<4->{red}?}
  \]
\end{frame}

%------------------------------------------------------------------------------%
\begin{frame}{Segunda Tentativa}
  Calculando
  \[
    F = \int \sqrt{a^2-x^2}\,dx
  \]
  \pause
  substituição
  \(
    \qquad
    x = a\sin(\theta)
    \qquad\qquad
    dx = a\cos(\theta)d\theta
    \vphantom{\frac{du}{2}}
  \)
  \pause
  \[
    F = \int \sqrt{a^2 - \big(a\sin(\theta)\big)^2\,} \,a\cos(\theta)d\theta
  \]
\end{frame}

%------------------------------------------------------------------------------%
\begin{frame}{Segunda Tentativa}
  \begin{align*}
    \onslide<+->{F & = \int \sqrt{a^2 - \big(a\sin(\theta)\big)^2\,} \,a\cos(\theta)d\theta     \\[2mm]}
    \onslide<+->{  & = \int \sqrt{a^2 \big( 1 - \sin^2(\theta)\big)\,} \, a\cos(\theta) d\theta \\[2mm]}
    \onslide<+->{  & = \int a^2 \sqrt{ \cos^2(\theta) } \,\cos(\theta) d\theta}
    \onslide<+->{  & & a > 0        \\[2mm]}
    \onslide<+->{  & = \int a^2 \cos^2(\theta) d\theta
                   & & \theta \in \left[-\frac{\pi}{2},\, \frac{\pi}{2}\right] \\[2mm]}
    \onslide<+->  {& = a^2 \int \cos^2(\theta) d\theta}
  \end{align*}
\end{frame}

%------------------------------------------------------------------------------%
\begin{frame}{Segunda Tentativa}
  \begin{align*}
    \onslide<+->{F & = a^2 \int \cos^2(\theta) d\theta         \\[2mm]}
    \onslide<+->{F & = \frac{a^2}{2} \int \big(1+\cos(2\theta)\big) d\theta
      & & \cos^2(\alpha)= \frac{1}{2}\big(1+\cos(2\alpha)\big) \\[2mm]}
    \onslide<+->{  & = \frac{a^2}{2}\left(\theta + \frac{\sin(2\theta)}{2}\right) + C \\[2mm]}
    \onslide<+->{  & = \frac{a^2}{2}\big(\theta + \sin(\theta)\cos(\theta)\big) + C \\[2mm]}
    \onslide<+->{  & = \frac{a^2}{2}
      \left[
      \arcsin\left(\frac{x}{a}\right)
      + \frac{x}{a}\sqrt{a^2-x^2}\,
      \right] + C
      & & \text{Justificativa em breve}}
  \end{align*}
\end{frame}

%------------------------------------------------------------------------------%
\begin{frame}{Transformação Inversa}
  \emph{Substituição simples} -- a nova variável é função da original
  \[
    u = f(x)
  \]

  \pause
  \emph{Substituição inversa} -- a variável original é função da nova
  \[
    x = g(u)
  \]
  \pause
  É necessário inverter $g$
\end{frame}

%------------------------------------------------------------------------------%
\begin{frame}{Substituições Trigonométricas}
\addtext{16}{0.35,2.2}{
\renewcommand*{\arraystretch}{2.3}
\begin{tabular}{c<{\;}>{\;}c<{\;}>{\;}c<{\;}>{\;}c} \hline
  Expressão
  & Substituição
  & Intervalo
  & Identidade \\ \hline
  \onslide<+->{{\color<4->{blue}\(\sqrt{a^2-x^2}\)}
  & {\color{blue}\(x=a\sin(\theta)\)}
  & \(-\frac{\pi}{2} \leq \theta \leq \frac{\pi}{2}\)
  & \(1-\sin^2(\theta) = \cos^2(\theta)\)} \\
  \onslide<+->{\(\sqrt{a^2+x^2}\)
  & \(x=a\tan(\theta)\)
  & \(-\frac{\pi}{2}< \theta < \frac{\pi}{2} \)
  & \(1+\tan^2(\theta)= \sec^2(\theta)\)} \\
  \onslide<+->{\(\sqrt{x^2-a^2}\)
  & \(x=a\sec(\theta)\)
  & \(0 \leq \theta < \frac{\pi}{2} \) ou \(\pi \leq \theta < \frac{3\pi}{2} \)
  & \(\sec^2(\theta)-1= \tan^2(\theta)\)} \\ \hline
\end{tabular}}
\end{frame}

%------------------------------------------------------------------------------%
%------------------------------------------------------------------------------%
\addtocounter{exemplo}{1}
\section{Exemplo \theexemplo}

%------------------------------------------------------------------------------%
\begin{frame}{Exemplo \theexemplo}
  Calcule \(\int \frac{\sqrt{9-x^2}}{x^2}dx\)

  \pause
  \vspace{0.5\baselineskip}
  Note
  \(
    \sqrt{9-x^2} = \sqrt{3^2-x^2}
  \)

  \pause
  \vspace{\baselineskip}
  Substituição inversa
  \[
    x = 3\sin(\theta)
    \qquad\qquad
    dx = 3 \cos(\theta) d\theta
    \qquad\text{com}\qquad
    -\frac{\pi}{2} \leq x \leq \frac{\pi}{2}
  \]
\end{frame}

%------------------------------------------------------------------------------%
\begin{frame}{Exemplo \theexemplo}
  \begin{align*}
    \onslide<+->{F & = \int \frac{\sqrt{9-x^2}}{x^2}dx          & & x = 3\sin(\theta)                              \\[2mm]}
    \onslide<+->{  & = \int \frac{\sqrt{9-9\sin^2(\theta)}}{9 \sin^2(\theta)} \, 3 \cos(\theta) d\theta            \\[2mm]}
    \onslide<+->{  & = \int \frac{3\times3}{9}\frac{\sqrt{1-\sin^2(\theta)}}{ \sin^2(\theta)} \cos(\theta) d\theta \\[2mm]}
    \onslide<+->{  & = \int \frac{\sqrt{1-\sin^2(\theta)}}{ \sin^2(\theta)} \cos(\theta) d\theta}
  \end{align*}
\end{frame}

%------------------------------------------------------------------------------%
\begin{frame}{Exemplo \theexemplo}
  \begin{align*}
    \onslide<+->{F & = \int \frac{\sqrt{1-\sin^2(\theta)}}{ \sin^2(\theta)} \cos(\theta) d\theta
                   & 1-\sin^2(\theta)= \cos^2(\theta) \\[2mm]}
    \onslide<+->{  & = \int \frac{\sqrt{\cos^2(\theta)}}{ \sin^2(\theta)} \cos(\theta) d\theta        \\[2mm]}
    \onslide<+->{  & = \int \frac{\cos^2(\theta)}{\sin^2(\theta)} d\theta                        \\[2mm]}
    \onslide<+->{  & = \int \cot^2(\theta) d\theta}
  \end{align*}
\end{frame}

%------------------------------------------------------------------------------%
\begin{frame}{Exemplo \theexemplo}
  \begin{align*}
    \onslide<+->{F & = \int \cot^2(\theta) d\theta       & \cot^2(\theta) = \csc^2(\theta)-1 \\[2mm]}
    \onslide<+->{  & = \int \csc^2(\theta) -1 \, d\theta & \big(\cot(\theta)\big)' = -\csc^2(\theta)\\[2mm]}
    \onslide<+->{  & = - \cot(\theta) - \theta +C        &}
  \end{align*}
  \onslide<+->{Precisamos voltar para a variável original \(x = 3\sin(\theta)\)}

  \onslide<+->{\[ \theta = \arcsin\left(\frac{x}{3}\right)\]}
  \onslide<+->{
  \[
    F = - {\color<7>{red}\cot\left(\arcsin\left(\frac{x}{3}\right)\right)}
        - \arcsin\left(\frac{x}{3}\right) + C
    \qquad
    \onslide<7>{\text{Forma ruim}}
  \]}
\end{frame}

%------------------------------------------------------------------------------%
\begin{frame}{Exemplo \theexemplo}
  \onslide<+->{\[
    x = 3\sin(\theta) \quad\Rightarrow\quad
    \sin(\theta)
    = \frac{x}{3}
    = \frac{\text{cateto oposto}}{\text{hipotenusa}}
  \]}

  \addgraphic{6}{10,3}{exemplo_triangulo_1}

  \onslide<+->{Por Pitágoras, o cateto adjacente é \(\sqrt{9-x^2}\)}

  \onslide<+->{Portanto
  \[
    \tan(\theta)
    = \frac{\text{cateto oposto}}{\text{cateto adjacente}}
  \]}

  \onslide<+->{\[
    \cot(\theta)
    = \frac{\text{cateto adjacente}}{\text{cateto oposto}}
    = \frac{\sqrt{9-x^2}}{x}
  \]}
\end{frame}

%------------------------------------------------------------------------------%
\begin{frame}{Exemplo \theexemplo}
  \begin{align*}
    \onslide<+->{\int \frac{\sqrt{9-x^2}}{x^2}dx
    & = -\cot(\theta) - \theta + C  \\[5mm]}
    \onslide<+->{& = -\frac{\sqrt{9-x^2}}{x} - \arcsin\left(\frac{x}{3}\right) + C}
  \end{align*}
\end{frame}

%------------------------------------------------------------------------------%
%------------------------------------------------------------------------------%
\addtocounter{exemplo}{1}
\section{Exemplo \theexemplo}

%------------------------------------------------------------------------------%
\begin{frame}{Exemplo \theexemplo}
  Calcule \(\int x \arcsin(x) dx\)

  \pause
  \vspace{\baselineskip}
  Vários métodos serão necessários
\end{frame}

%------------------------------------------------------------------------------%
\begin{frame}{Exemplo \theexemplo\ -- Integral por Partes}
  \onslide<+->{\[
    F = \int x \arcsin(x) dx
    \]}

  \onslide<+->{Integral por partes}
  \begin{align*}
    \onslide<+->{u  & = \arcsin(x)              & dv & = xdx           \\}
    \onslide<+->{du & =\frac{1}{\sqrt{1-x^2}}dx & v  & = \frac{x^2}{2}}
  \end{align*}

  \[
  \onslide<+->{F = uv - \int v du}
  \onslide<+->{  = \frac{x^2}{2} \arcsin(x)}
  \onslide<+->{  - \frac{1}{2}\int \frac{x^2}{\sqrt{1-x^2}}dx}
  \]
\end{frame}

%------------------------------------------------------------------------------%
\begin{frame}{Exemplo \theexemplo\ -- Substituição Trigonométrica}
  \onslide<+->{Resolver \( \int \frac{x^2}{\sqrt{1-x^2}}dx \)}

  \vspace{\baselineskip}
  \onslide<+->{Substituição inversa
    \[
    x  = \sin(\theta)
    \qquad
    dx = \cos(\theta)d\theta
    \quad\text{com}\quad
    -\frac{\pi}{2} \leq \theta \leq \frac{\pi}{2}
    \]}
  \[
  \onslide<+->{G = \int \frac{x^2}{\sqrt{1-x^2}}dx}
  \onslide<+->{  = \int \frac{\sin^2(\theta)}{\sqrt{1-\sin^2(\theta)}} \cos(\theta)d\theta}
  \]
\end{frame}

%------------------------------------------------------------------------------%
\begin{frame}{Exemplo \theexemplo\ -- Substituição Trigonométrica}
  \begin{align*}
    \onslide<+->{G & = \int \frac{x^2}{\sqrt{1-x^2}}dx                                         \\[2mm]}
    \onslide<+->{  & = \int \frac{\sin^2(\theta)}{\sqrt{1-\sin^2(\theta)}} \cos(\theta)d\theta \\[2mm]}
    \onslide<+->{  & = \int \frac{\sin^2(\theta)}{\sqrt{\cos^2(\theta)}} \cos(\theta)d\theta   \\[2mm]}
    \onslide<+->{  & = \int \sin^2(\theta) d\theta}
  \end{align*}
\end{frame}

%------------------------------------------------------------------------------%
\begin{frame}{Exemplo \theexemplo\ -- Integral Trigonométrica}
  \onslide<+->{Calcular \( G = \int \sin^2(\theta) d\theta \)}

  \vspace{\baselineskip}
  \onslide<+->{Manipulação algébrica
    \[
    \sin^2(\theta)= \frac{1}{2}\big(1-\cos(2 \theta)\big)
    \]}
\end{frame}

%------------------------------------------------------------------------------%
\begin{frame}{Exemplo \theexemplo\ -- Integral Trigonométrica}
  \begin{align*}
    \onslide<+->{G & = \int \sin^2(\theta) d \theta                                            \\[2mm]}
    \onslide<+->{  & = \frac{1}{2}\int\big(1-\cos(2 \theta)\big) d\theta                       \\[2mm]}
    \onslide<+->{  & = \frac{1}{2}\left(\theta -\frac{\sin(2\theta)}{2}\right) + C             \\[2mm]}
    \onslide<+->{  & = \frac{1}{2}\left(\theta -\frac{2\sin(\theta)\cos(\theta)}{2}\right) + C \\[2mm]}
    \onslide<+->{  & = \frac{1}{2}\big(\theta -\sin(\theta)\cos(\theta)\big) + C}
  \end{align*}
\end{frame}

%------------------------------------------------------------------------------%
\begin{frame}[t]{Exemplo \theexemplo\ -- Desfazendo a Transformação Trigonométrica}
  \[
  \onslide<+->{ x
    = \sin(\theta)
    = \frac{\text{cateto oposto}}{\text{hipotenusa}}}
  \qquad
  \onslide<+->{
    \theta=\arcsin(x)}
  \]

  \addgraphic{4}{10,3}{exemplo_triangulo_4}

  \onslide<+->{\[
    \cos(\theta)= \sqrt{1-x^2}
    \]}
  \begin{align*}
    \onslide<+->{G & = \int \frac{x^2}{\sqrt{1-x^2}}dx                       \\[2mm]}
    \onslide<+->{  & = \frac{1}{2}\big(\theta -\sin(\theta)\cos(\theta)\big) \\[2mm]}
    \onslide<+->{  & = \frac{1}{2}\left( \arcsin(x)-x\sqrt{1-x^2} \right) + C}
  \end{align*}
\end{frame}

%------------------------------------------------------------------------------%
\begin{frame}{Exemplo \theexemplo\ -- Retornando a Integral Original}
  \begin{align*}
    \onslide<+->{F & = \int x \arcsin(x) dx                                 \\[5mm]}
    \onslide<+->{  & = \frac{x^2}{2} \arcsin(x)
      - \frac{1}{2}\int \frac{x^2}{\sqrt{1-x^2}}dx           \\[5mm]}
    \onslide<+->{  & = \frac{x^2}{2} \arcsin(x)
      - \frac{1}{4}\left(\arcsin(x)-x\sqrt{1-x^2}\right) + C \\[5mm]}
    \onslide<+->{  & = \left(\frac{x^2}{2} - \frac{1}{4}\right) \arcsin(x)
      + \frac{x}{4} \sqrt{1-x^2}
      + C}
  \end{align*}
\end{frame}

%------------------------------------------------------------------------------%
\section{Lista Mínima}

%------------------------------------------------------------------------------%
\begin{frame}{Lista Mínima}
  Estudar a Seção 4.4 da Apostila

  Exercícios:

  \vfill
  Atenção: A prova é baseada no livro, não nas apresentações
\end{frame}

%------------------------------------------------------------------------------%
\end{document}
%------------------------------------------------------------------------------%